% Part: history
% Chapter: biographies 
% Section: alonzo-church

\documentclass[../../../include/open-logic-section]{subfiles}

\begin{document}

\olfileid{his}{bio}{chu}

\olsection{అలోంజో చర్చ్}

\olphoto{church-alonzo}{అలోంజో చర్చ్}

అలోంజో చర్చ్ 1903 జూన్ 14న వాషింగ్టన్, డి.సి.లో జన్మించారు. చిన్నతనంలో గాలితో కాల్చే తుపాకీకి సంబంధించిన ప్రమాదం వల్ల ఆయన ఒక కంటిలో చూపును కోల్పోయారు. 1920లో కనెక్టికట్‌లో సన్నాహక పాఠశాల విద్యను ముగించి, అదే సంవత్సరంలో ప్రిన్స్‌టన్‌లో విశ్వవిద్యాలయ చదువు మొదలుపెట్టారు. 1927లో డాక్టరేట్ పూర్తిచేశారు. కొన్ని సంవత్సరాలు విదేశాల్లో గడిపిన తరువాత చర్చ్ ప్రిన్స్‌టన్‌కు తిరిగి వచ్చారు. ఆయన అత్యంత మర్యాదగా, జాగ్రత్తగా వ్యవహరించేవారని పేరుంది. బ్లాక్‌బోర్డుపై ఆయన రాత ఎంతో శుభ్రంగా ఉండేది; ముఖ్యమైన పత్రాలను పారదర్శక జిగురైన డూకో సిమెంట్‌తో జాగ్రత్తగా పూతపూసి భద్రపరిచేవారు. విద్యా వ్యాసంగం వెలుపల ఆయనకు విజ్ఞాన కల్పన పత్రికలు చదవడం ఇష్టం; రచనలో తప్పులు కనిపిస్తే సంపాదకులకు రాయడానికి వెనుకాడేవారు కాదు.

చర్చ్ విద్యాసంబంధ సాధనలు విశేషమైనవి. తన శిష్యులు స్టీఫెన్ క్లీన్, బార్క్లీ రోసర్‌లతో కలిసి, అలన్ ట్యూరింగ్ ట్యూరింగ్ యంత్రాన్ని అభివృద్ధి చేయడానికి స్వతంత్రంగా, ప్రభావకంగా గణించగలగడంపై లాంబ్డా కలనశాస్త్రం అనే సిద్ధాంతాన్ని రూపొందించారు. గణనీయతకు సంబంధించిన ఈ రెండు నిర్వచనాలు తుల్యమైనవి. వాటి నుంచి నేడు \emph{చర్చ్--ట్యూరింగ్ సిద్ధాంతప్రతిపాదన} అని పిలిచే వాదన వస్తుంది: సహజ సంఖ్యలపై ఒక ప్రమేయం ప్రభావకంగా గణించదగినది అప్పుడూ అప్పుడే అది ట్యూరింగ్ యంత్రంతో (లేదా లాంబ్డా కలనశాస్త్రంతో) గణించదగినది. నేడు \emph{చర్చ్ సిద్ధాంతం} అని పిలిచే ఫలితాన్నీ ఆయన నిరూపించారు: ప్రథమ శ్రేణి సూత్రాల చెల్లుబాటును నిర్ణయించే సమస్య అనిర్ణయనీయమైనది.

చర్చ్ వృద్ధాప్యంలోనూ తన కృషిని కొనసాగించారు. 1967లో ప్రిన్స్‌టన్‌ను వదిలి యూసీఎల్‌ఏకు వెళ్లారు; 1990లో పదవీ విరమణ చేసే వరకు అక్కడ ఆచార్యునిగా ఉన్నారు. 1995 ఆగస్టు 1న, 92 ఏళ్ల వయసులో చర్చ్ మరణించారు.

\begin{reading}
చర్చ్ సంక్షిప్త జీవిత చరిత్రకు \citet{EndertonND} చూడండి. లాంబ్డా కలనశాస్త్రం, నిర్ణయ సమస్య (చర్చ్ సిద్ధాంతప్రతిపాదన) గురించి చర్చ్ రాసిన మూల వ్యాసాలు \citet{Church1936,Church1936a}. 1930ల ప్రిన్స్‌టన్ గణిత సమాజం గురించి చర్చ్‌తో జరిగిన ముఖాముఖిని \citet{Aspray1984} నమోదు చేశారు. చర్చ్ 1936 నుంచి 1979 వరకు \emph{Journal of
Symbolic Logic}లో పుస్తక సమీక్షల పరంపర రాశారు. అవన్నీ జాన్ మాక్‌ఫార్లేన్ వెబ్‌సైట్‌లో భద్రపరచబడ్డాయి \citep{MacFarlane2015}.
\end{reading}
\end{document}
